% RESEARCH_DRAFT_SHA256: b7782871ce53d0a754ec53afd4ad23b1ce36153fe8ef912263b43f8a0eeb9974
% Layout and source-reference wrappers expanded for the web renderer.
\documentclass[11pt]{article}
\usepackage[margin=25mm]{geometry}
\usepackage{fontspec,unicode-math}
\setmainfont{Latin Modern Roman}
\setmathfont{Latin Modern Math}
\usepackage{amsmath,amssymb,mathtools,xurl,hyperref}
\hypersetup{colorlinks=true,urlcolor=blue,linkcolor=blue}
\setcounter{secnumdepth}{0}
\setlength{\parindent}{0pt}
\setlength{\parskip}{5pt}
\setlength{\emergencystretch}{3em}
\newcommand{\sref}[2]{\hyperlink{source:#1:#2}{[S#2]}}
\newcommand{\cataloguescope}[1]{\paragraph{Source record.}#1}
\newcommand{\floor}[1]{\left\lfloor#1\right\rfloor}
\begin{document}
\section*{Research-attempt provenance}
This standalone research attempt retains the complete original definition body below,
including its historical source and preparation notes. Those notes describe the
earlier catalogue preparation; the new attempt and its verification are identified
separately after that body. No original ID, catalogue status, or source statement is changed.

% UnsolvedMath ID 30003755; source catalogue code OWR-15999-001
\section{Unified Simplicity Models for Massive Inverse Problems}\label{30003755}
{\small UnsolvedMath ID 30003755 · Computer Science}
\subsection{Definitions and mathematical statement}

\paragraph{Shared notation.}Standard number systems are $\mathbb N,\mathbb Z,\mathbb Q,\mathbb R,\mathbb C$. Specialized notation and index ranges follow the cited source and the definitions below.


A sparse object has few nonzero coordinates; a low-rank matrix has a low-dimensional column space. Entropy quantifies uncertainty of a probability distribution. These are distinct ways to impose a reduced-description structure on an inverse problem.
\paragraph{Record scope and qualifications.}
The catalogue title is retained as metadata. Named certificates, contracts, bounds or frameworks in that title are not asserted as results of the original source extract.
The requested unification is a mathematical modeling direction rather than a quantified equivalence or reconstruction theorem.

Source-record limitation: A unified notion of simplicity is requested without a specific mathematical criterion.
\paragraph{Statement recorded in the supplied source.}
In phase retrieval and inverse problems for massive complex data, determine how sparsity and minimal information complexity can be used, and whether a unified notion of simplicity can include sparsity, low rank, and low entropy as special cases. \href{https://doi.org/10.4171/owr/2018/14}{[S2]}
\paragraph{Source status.}
The supplied catalogue labels this source record open; that metadata was checked against the supplied snapshot on 2026-10-11. The label is preserved for provenance and is not an independent verification of present open status or mathematical correctness. The catalogue literature assessment is dated 2026-08-21.
\subsection{Short English statement}
Seek a unified simplicity concept for large inverse problems.
\subsection{Sources}
\begin{itemize}
\item[S1] Ulam AI and the UnsolvedMath Contributors, UnsolvedMath: A Curated Collection of Open Mathematics Problems, record 30003755 (OWR-15999-001). Catalogue curation and metadata: CC BY 4.0. The source extract is retained; definitions, English summaries and qualifications are editorial additions.
\url{https://huggingface.co/datasets/ulamai/UnsolvedMath}.
\item[S2] Applied Harmonic Analysis and Data Processing (2019), Oberwolfach Reports, DOI 10.4171/owr/2018/14, pp. 723--792
\url{https://doi.org/10.4171/owr/2018/14}.
\item[S3] Oberwolfach Report 14
\url{https://doi.org/10.4171/OWR/2018/14}.
\item[S4] Sparse phase retrieval
\url{https://doi.org/10.1016/j.acha.2014.04.001}.
\item Source-specific Oberwolfach Reports licences, including CC BY-SA 4.0 where applicable, are retained; see THIRD\_PARTY\_NOTICES.md.
\end{itemize}
\paragraph{Preparation and review.}
The original source and UnsolvedMath contributors retain attribution for the source extract. Codex (GPT-6; exact serving revision not exposed) assisted with definitions, English summaries, scope notes and formatting. No new human mathematical authorship or independently verified solution is asserted. Permanent source ID 30003755 and code OWR-15999-001 checked against the complete supplied registry and current repository. Distinct upstream records, including repeated questions and broad or damaged extracts, are retained for complete coverage of the source CS category. A repeated record is not a new independent conjecture. Complete mathematical-equivalence and current-literature review has not been performed. Preparation and final review follow contribution prompts C and D, with source gaps disclosed. Structural and rendering checks do not establish a new theorem or open status.
\label{end:30003755}

\section{Literature and context check}
The introduction to OWR 14/2018, p.~724, was read. It poses the proposed unification of sparsity, low rank and low entropy as a broad modeling direction, alongside nonlinear phase retrieval. It does not supply a formal universal definition or a desired equivalence theorem. The framework below is an elementary conditional identifiability criterion, followed by examples showing why one common convex penalty cannot simply be asserted to cover all three notions.
\section{Research attempt: Secant geometry unifies identifiability but not all simplicity penalties}

\paragraph{A precise common recovery criterion.}
Let $\mathcal M$ be a specified nonempty subset of a finite-dimensional normed vector space and let $A$ be a linear measurement map. Exact recovery from $Ax$ for every $x\in\mathcal M$ is possible if and only if
\[
 \ker A\cap(\mathcal M-\mathcal M)=\{0\}.
\]
Indeed two model points have the same measurement exactly when their difference lies in this intersection. If the intersection is trivial, assigning to each measurement in $A\mathcal M$ its unique preimage gives a set-theoretic recovery map; this asserts neither efficient computation nor numerical stability.

A stronger condition quantifies stability: suppose $\|A(x-z)\|\ge\alpha\|x-z\|$ for all $x,z\in\mathcal M$, with $\alpha>0$. If $y=Ax+e$, $\|e\|\le\epsilon$, and a returned model point $\widehat x$ satisfies $\|A\widehat x-y\|\le\epsilon$, then
\[
 \|\widehat x-x\|\le2\epsilon/\alpha.
\]
This follows by the triangle inequality and the displayed lower bound. An algorithm that finds such a point is an additional requirement, not part of the geometric criterion.

\paragraph{Sparse vectors and low-rank matrices fit this criterion.}
For the set of vectors with at most $k$ nonzero coordinates, its difference set is exactly the set of vectors with at most $2k$ nonzeros, with the ambient dimension capping that count. One inclusion follows by taking the union of supports. For the other, split a support of size at most $2k$ into two sets of size at most $k$, placing one part in $x$ and the negative of the other in $z$. The difference is the desired vector.

For real matrices of rank at most $r$, the difference set is exactly the matrices of rank at most $2r$, again capped by dimensions. The upper bound follows because the column space of a difference lies in the sum of the two column spaces. Conversely, a singular value decomposition of a matrix of rank at most $2r$ can be split into two sums with at most $r$ terms, putting the negative of the second sum in $z$. The Frobenius norm then gives the same stability criterion. These parallel facts explain a common measurement-geometric language; they do not assert that support size and rank are the same intrinsic quantity or that the corresponding efficient algorithms coincide.

\paragraph{Entropy exposes a failed convexification.}
For probability vectors $p\in\mathbb R^N$, let $H(p)=-\sum_i p_i\log p_i$, with $0\log0=0$. For every $h\ge0$, the set $\mathcal E_h=\{p:H(p)\le h\}$ contains every point mass $e_i$. Its convex hull is therefore the entire probability simplex, since $p=\sum_i p_i e_i$ for every probability vector. Consequently replacing a low-entropy model by its convex hull erases the entropy restriction completely. An $\ell_1$ penalty is also constant on this simplex and cannot restore the missing distinction.

Low entropy does not imply exact sparsity either. For $0<\epsilon<1$,
\[
 p_\epsilon=\left(1-\epsilon,\frac{\epsilon}{N-1},\ldots,
                              \frac{\epsilon}{N-1}\right)
\]
has full support for $N\ge2$, while
\[
 H(p_\epsilon)=-(1-\epsilon)\log(1-\epsilon)-\epsilon\log\epsilon
                         +\epsilon\log(N-1)\longrightarrow0
\]
as $\epsilon\downarrow0$. In the other direction a distribution on at most $k$ points has $H(p)\le\log k$: writing its positive support size as $s\le k$, the inequality $\log t\le t-1$ gives $\sum_i p_i\log((1/s)/p_i)\le\sum_i(1/s-p_i)=0$. Thus entropy and support are related but cannot be interchanged as exact model constraints.

\paragraph{Phase retrieval requires a quotient and a scale condition.}
For real intensity measurements $\mathcal A(x)=(|a_i^Tx|^2)_i$, one always has $\mathcal A(x)=\mathcal A(-x)$. Exact vector identifiability on a model containing both signs is impossible; the target must be a sign-equivalence class if that is the intended information. Even after quotienting, a global inverse-Lipschitz estimate down to zero cannot hold for a nontrivial cone: $\|\mathcal A(tx)-\mathcal A(0)\|=t^2\|\mathcal A(x)\|$ while the distance from the sign class of $tx$ to zero is $t\|x\|$. Their ratio tends to zero. A norm normalization or a weaker stability modulus is needed. Merely declaring an object sparse, low rank or low entropy does not fix these measurement-induced ambiguities.
\section{Verification and adversarial audit}
Both directions of the secant criterion and both difference-set identities are proved. The stable recovery bound is conditional on a feasible returned model point and does not hide an optimization algorithm. The entropy convex-hull failure, full-support small-entropy family, support-to-entropy inequality and phase scaling counterexample were checked including the zero vector and sign symmetry.
\paragraph{Audit corrections.} Distinguished set-theoretic identifiability, stability and efficient recovery. Rejected equating low entropy with sparsity or taking its convex hull as a faithful penalty. Added the sign and small-amplitude obstructions for phase retrieval.
\paragraph{External computation.} None; symbolic arguments and boundary checks only.
\section{Scope and remaining gaps}
The original record requests a unifying modeling theory without a fixed success criterion. The conditional secant language gives a common identifiability statement, but no universal simplicity functional, efficient reconstruction method or full nonlinear theory is established. Low-entropy convexification and intensity ambiguities are concrete obstacles to the proposed naive unification, not a refutation of every possible framework.
\paragraph{Integration review.} Codex (GPT-6; exact serving revision not exposed) reviewed the written arguments and expanded layout and source-reference wrappers. Repository metadata and the production web renderer were checked using Python 3.13 and JavaScript with MathJax 3.2.2. This model-assisted review is not independent human proof verification; no PDF compilation or formal proof checking is asserted.

\section{Completion Estimate}
\noindent\textbf{COMPLETION: 0\%}
The calculations below the preserved source statement investigate a concrete route,
but establish no substantiated new progress toward the entire original target.
Known special cases and failed sufficient conditions are not credited as a solution.
\section{References}
The original source references and catalogue attribution are preserved in the
Sources section above. The following URLs were actually consulted for this attempt;
retention of an additional original reference does not claim fresh consultation of it.
\begin{itemize}
\item \url{https://ems.press/content/serial-article-files/46736} (consulted 2026-10-11). 
\end{itemize}
\section{Model and contributor disclosure}
OpenAI gpt-6-astra, requested reasoning configuration ultra, generated this research
attempt on 2026-10-11. The exact serving revision is not exposed. Codex GPT-6
performed orchestration, repository-format review, and written-proof review of this attempt; its exact serving revision
is also not exposed. The mathematical route, boundary checks and final adversarial
self-audit followed contribution prompts A then B in
\texttt{docs/contribution\_prompts.md}, after reading \texttt{CONTRIBUTING.md}.
Python 3.14.2 assembled the standalone documents; any mathematical computation is
specified in Verification. No theorem prover was used. No human mathematical
authorship, substantial human proof correction, or independent human proof checking
is asserted. Model-assisted review, source retrieval, and structural checks do not
establish correctness or novelty. Full private reasoning transcripts are unavailable
as reproducibility artifacts; the complete arguments, inputs and outcomes actually
used are stated here or in the companion batch audit.
\section{Final status}
Unresolved as a new research contribution. This is a concrete route investigation,
not a claim of new complete resolution or a current-status certification. A known
answer reproduced for context does not become a new contribution. The catalogue's
source-status label is preserved; this attempt status does not assert that a
source-labelled solved question is open.
\end{document}
